% Guia de observacion no participante

\documentclass[11pt,a4paper]{article}

\usepackage[utf8]{inputenc}
\usepackage[T1]{fontenc}
\usepackage[spanish]{babel}
\usepackage[scaled=0.92]{helvet}
\renewcommand{\familydefault}{\sfdefault}

% interlineado general
\linespread{1.2}
\usepackage[a4paper,top=2.4cm,bottom=2.3cm,left=2.3cm,right=2.3cm,headheight=14pt,headsep=0.4cm]{geometry}
\usepackage[table]{xcolor}
\usepackage{enumitem}
\usepackage{fancyhdr}
\usepackage{lastpage}
\usepackage{booktabs}
\usepackage{tabularx}
\usepackage{longtable}
\usepackage{needspace}
\usepackage{array}
\usepackage{amssymb}
\usepackage[most]{tcolorbox}
\usepackage[hidelinks]{hyperref}

% Colores institucionales
\definecolor{uteqverde}{HTML}{2C5F2D}
\definecolor{uteqazul}{HTML}{1E4D8B}
\definecolor{uteqoro}{HTML}{D4A017}
\definecolor{rowgray}{HTML}{F2F4F2}

% Encabezado y pie
\pagestyle{fancy}
\fancyhf{}
\fancyfoot[L]{\scriptsize\color{uteqverde} ISR-401 Ingeniería de Requerimientos $\cdot$ Guía de observación}
\fancyfoot[R]{\scriptsize\color{uteqverde} Página \thepage\ de \pageref{LastPage}}
\renewcommand{\headrulewidth}{0pt}
\renewcommand{\footrulewidth}{0.4pt}
\renewcommand{\footrule}{{\color{uteqoro}\hrule width\headwidth height\footrulewidth \vskip\footruleskip}}

% Estilo de seccion
\newcommand{\seccion}[1]{%
	\par\vspace{16pt}%
	\noindent{\color{uteqverde}\large\bfseries #1}\par%
	\vspace{-3pt}{\color{uteqoro}\rule{\linewidth}{1pt}}\par\vspace{6pt}%
}

\newcommand{\etiqueta}[1]{\textcolor{uteqverde}{\bfseries #1}}
\newcommand{\cb}{$\square$}

% titulo de bloque: reserva espacio para que no quede solo al final de la pagina
\newcommand{\bloque}[2]{\Needspace{#1}\par\textcolor{uteqazul}{\bfseries #2}\par\vspace{8pt}}
\newcommand{\encab}[1]{\textcolor{white}{\textbf{#1}}}
\newcolumntype{Y}{>{\raggedright\arraybackslash}X}
\newcolumntype{R}[1]{>{\raggedright\arraybackslash}p{#1}}
\newcolumntype{C}[1]{>{\centering\arraybackslash}p{#1}}

% longtable pegada a los margenes y sin espacio extra antes y despues
\setlength{\LTpre}{0pt}
\setlength{\LTpost}{0pt}
\setlength{\LTleft}{0pt}
\setlength{\LTright}{0pt}
\setlength{\parindent}{0pt}
\setlength{\parskip}{6pt}
\setlist[itemize,1]{label=\textcolor{uteqoro}{\textbullet},itemsep=4pt,topsep=4pt,leftmargin=16pt}

% Lista de cotejo: codigo, aspecto, Si, Parcial, No, registro
\newenvironment{cotejo}{%
	\small\setlength{\parskip}{0pt}\renewcommand{\arraystretch}{2.0}%
	\rowcolors{2}{white}{rowgray}%
	\longtable{@{}l R{5.8cm} C{0.7cm} C{1.0cm} C{0.7cm} R{5.0cm}@{}}%
		\rowcolor{uteqverde}%
		\encab{Cód.} & \encab{Qué se observa} & \encab{Sí} & \encab{Parc.} & \encab{No} & \encab{Hecho observado} \\
		\endhead
	}{\endlongtable\par\vspace{10pt}}

\newcommand{\fila}[2]{#1 & #2 & \cb & \cb & \cb & \\}

\begin{document}

	\begin{center}
		{\small\color{uteqazul}\bfseries FACULTAD DE CIENCIAS DE LA COMPUTACIÓN $\cdot$ CARRERA DE SOFTWARE}\\[6pt]
		{\LARGE\color{uteqverde}\bfseries GUÍA DE OBSERVACIÓN}\\[3pt]
		{\normalsize Proyecto Fin de Curso $\cdot$ Ingeniería de Requerimientos (ISR-401)}
	\end{center}

	\begin{tcolorbox}[enhanced,colback=rowgray,colframe=uteqverde,boxrule=0.6pt,arc=1.5mm,left=4pt,right=4pt,top=3pt,bottom=3pt]
		\small
		\etiqueta{Objetivo:} Observar, sin intervenir (observación no participante), cómo se realizan actualmente en el local el seguimiento de los pedidos personalizados, las ventas, el cobro de abonos y las entregas, qué herramientas se usan y qué dificultades aparecen.
	\end{tcolorbox}

	\seccion{1. Datos de la sesión}

	\begingroup
	\small
	\renewcommand{\arraystretch}{2.0}
	\begin{tabularx}{\textwidth}{@{}p{5cm} Y@{}}
		\toprule
		Código de la sesión & \makebox[4cm]{\hrulefill} \\
		Lugar & \makebox[8cm]{\hrulefill} \\
		Hora de inicio y de fin & \makebox[3cm]{\hrulefill} \ a \ \makebox[3cm]{\hrulefill} \\
		Persona observada (rol y código) & \makebox[4cm]{\hrulefill} \ $\cdot$ \ \makebox[3cm]{\hrulefill} \\
		Consentimiento informado firmado & \cb\ Sí \quad \cb\ No (si es No, no se inicia la observación) \\
		Autorización para observar el flujo de chats (sin anotar su contenido) & \cb\ Sí \quad \cb\ No \\
		\bottomrule
	\end{tabularx}
	\endgroup

	\seccion{2. Reglas de la observación}

	\begin{itemize}
		\item El observador no interviene, no corrige ni sugiere soluciones durante la sesión; solo registra.
		\item Se anotan hechos (qué hizo, con qué herramienta, cuánto tardó, qué falló), no opiniones. Toda interpretación se escribe aparte con la etiqueta \emph{[interpretación]}.
		\item No se anotan nombres, teléfonos, direcciones, cédulas, datos bancarios ni contenido de conversaciones. Los clientes se registran con un código (C1, C2, \ldots) y los pedidos como P1, P2, \ldots
		\item No se toman fotografías ni videos salvo autorización escrita y sin rostros ni datos identificables.
		\item Si la persona observada pide pausar o detener la observación, se hace de inmediato y se deja constancia.
	\end{itemize}

	\Needspace{5cm}
	\seccion{3. Lista de cotejo}

	Marcar \etiqueta{Sí} (se observó completo), \etiqueta{Parc.} (parcial o con dificultades) o \etiqueta{No} (no ocurrió o no se pudo observar). En la última columna se escribe el hecho observado y la herramienta usada (cuaderno, WhatsApp, hoja de cálculo, otra).

	\vspace{4pt}

	\bloque{9.5cm}{A. Pedidos personalizados}

	\begin{cotejo}
		\fila{A-01}{Por qué canal llega el pedido y cómo se registra (cuaderno, mensaje, memoria, otro).}
		\fila{A-02}{Qué características del pedido se anotan y cómo se manejan los cambios posteriores.}
		\fila{A-03}{Cómo se consulta el estado de un pedido y cómo se informa al cliente del avance o los retrasos.}
		\fila{A-04}{Se pierden, repiten o se contradicen datos de un pedido.}
	\end{cotejo}

	\bloque{7.5cm}{B. Ventas de productos terminados}

	\begin{cotejo}
		\fila{B-01}{Cómo se verifica qué productos hay disponibles antes de ofrecerlos.}
		\fila{B-02}{Cómo se registra una venta (si se registra).}
		\fila{B-03}{Algún producto solicitado que no estaba disponible o cuya existencia no se pudo confirmar.}
	\end{cotejo}

	\bloque{9cm}{C. Pagos y abonos}

	\begin{cotejo}
		\fila{C-01}{Formas de pago y cómo se registra un abono o un pago total.}
		\fila{C-02}{Cómo se verifica una transferencia.}
		\fila{C-03}{Cómo se calcula el saldo pendiente y cómo se confirma el abono al cliente.}
		\fila{C-04}{Alguna diferencia, duda o demora al verificar lo cobrado.}
	\end{cotejo}

	\bloque{7.5cm}{D. Entregas}

	\begin{cotejo}
		\fila{D-01}{Tipo de entrega (en el local, a domicilio, encomienda) y qué datos de destino se piden.}
		\fila{D-02}{Cómo se confirma que el cliente recibió el pedido y cómo se cierra el pedido.}
		\fila{D-03}{Algún problema en la entrega (retraso, dato faltante, confirmación pendiente).}
	\end{cotejo}

	\Needspace{9.5cm}
	\seccion{4. Registro cronológico de eventos}

	Se anotan solo los eventos relevantes, en orden, con códigos de cliente y de pedido.

	\vspace{4pt}

	\begingroup
	\small
	\setlength{\parskip}{0pt}
	\renewcommand{\arraystretch}{1.9}
	\rowcolors{2}{white}{rowgray}
	\begin{longtable}{@{}C{1.3cm} R{4.8cm} C{1.8cm} C{2.3cm} R{4.5cm}@{}}
		\rowcolor{uteqverde}
		\encab{Hora} & \encab{Hecho observado} & \encab{Cód. (A--D)} & \encab{Herramienta} & \encab{[Interpretación]} \\
		\endhead
		& & & & \\
		& & & & \\
		& & & & \\
		& & & & \\
		& & & & \\
	\end{longtable}
	\endgroup

	\Needspace{6.5cm}
	\seccion{5. Dificultades observadas}

	\begingroup
	\small
	\setlength{\parskip}{0pt}
	\renewcommand{\arraystretch}{2.0}
	\rowcolors{2}{white}{rowgray}
	\begin{longtable}{@{}C{1.1cm} R{4.6cm} R{4.7cm} R{4.7cm}@{}}
		\rowcolor{uteqverde}
		\encab{N.º} & \encab{Situación observada} & \encab{Efecto en el trabajo (tiempo perdido, error, reproceso)} & \encab{Duda por aclarar} \\
		\endhead
		1 & & & \\
		2 & & & \\
		3 & & & \\
		4 & & & \\
	\end{longtable}
	\endgroup

	\seccion{6. Verificación y cierre}

	\begin{tcolorbox}[enhanced,colback=rowgray,colframe=uteqverde,boxrule=0.6pt,arc=1.5mm,left=4pt,right=4pt,top=3pt,bottom=3pt]
		\small
		\cb\ Las notas no contienen nombres, teléfonos, direcciones, cédulas, datos bancarios ni contenido de chats.\\
		\cb\ Los clientes y pedidos se registraron solo con código.\\
		\cb\ Los hechos observados están separados de las interpretaciones.\\
		\cb\ Las fotografías, si las hubo, contaron con autorización escrita y no muestran rostros ni documentos legibles.
	\end{tcolorbox}

	\vspace{14pt}
	\begin{tabularx}{\textwidth}{@{}Y Y@{}}
		Observador/a: \makebox[4.5cm]{\hrulefill} & Fecha: \makebox[3.5cm]{\hrulefill} \\
	\end{tabularx}

\end{document}